\section{Introduction}

% Seeded from DOCX text extraction (see scratch/by_section/10_introduction.tex).
% TODO: Verify wording + fill in missing equations from the PDF.

For nearly a century, gravitational anomalies at galactic and cosmological scales have been parameterized by an additional non-luminous sector. The $\Lambda$CDM concordance model fits the CMB, BAO, and large-scale structure with high precision when cold dark matter (CDM) is treated as an effectively pressureless gravitating component.

Yet the microphysical identity of that dominant non-baryonic component remains empirically unresolved. Many canonical weak-scale interaction scenarios have been constrained by direct detection, indirect searches, and collider experiments. This motivates renewed attention to alternatives that treat the dark sector first as an empirical role and only secondarily as an ontological identity.

This paper develops an inverse-GR alternative consistent with the Discussion framing in \S9: we retain standard GR on the left-hand side and ask what effective stress--energy must appear on the right-hand side if we do not assume particulate CDM. The organizing equation is \emph{(see PDF)}, where \emph{(see PDF)} is defined as the effective response sector required by the observed phenomenology. This definition is ontologically agnostic; we do not claim spacetime is literally a fluid or solid, only that it admits an EFT/constitutive response closure.

In the circular-orbit regime, the data provide \emph{(see PDF)} and therefore \emph{(see PDF)}. The baryonic prediction \emph{(see PDF)} is assembled from SPARC's stellar and gas components with specified mass-to-light ratios. The model-independent residual \emph{(see PDF)} defines the ``missing curvature'' attributable in $\Lambda$CDM to a halo, and attributable here to the response sector.

The inverse mapping alone is underdetermined: many effective sources could reproduce a given \emph{(see PDF)}. The scientific content of this work is therefore compression. We impose a boundary-layer closure in which the response activates near a baryonic transition radius \emph{(see PDF)} (operationally where \emph{(see PDF)}) and satisfies an approximate flux-conservation law outside the activation zone, yielding an asymptotic \emph{(see PDF)} tail with $\le 2$ degrees of freedom per galaxy.

We further distinguish \emph{(see PDF)} (a global fit amplitude under the imposed closure) from \emph{(see PDF)} (a robust outer estimator computed directly from the outer \emph{(see PDF)} subset). Their strong agreement in rank across SPARC-175 supports the interpretation that the inferred response amplitude tracks genuine outer physics, while deviations diagnose where a single fixed shape requires extension.

Rotation curves primarily constrain the Newtonian-gauge potential \emph{(see PDF)}; lensing constrains the Weyl potential \emph{(see PDF)}. Consistent with \S9, we treat lensing as a gateway discriminant between Branch A (metric equivalence, \emph{(see PDF)}) and Branch B (slip-allowed, \emph{(see PDF)} localized to activation zones).

\begin{figure}[t]
	\centering
	% TODO: Add schematic figure file under figures/ and include it here.
	% \includegraphics[width=0.9\linewidth]{figures/figure1_concept.pdf}
	\caption{Conceptual schematic of the response-sector framework. The central luminous disk represents the baryonic mass distribution. Concentric lobes illustrate the radial residual $\Delta a$, with compression ($\Delta a > 0$, orange) and rarefaction ($\Delta a < 0$, blue) zones. The dashed ellipse marks the transition radius $R_t$ where \emph{(see PDF)}. Two amplitude estimators are distinguished: \emph{(see PDF)} (robust outer-subset Huber location) and \emph{(see PDF)} (single-parameter model-closure fit). The oscillatory structure is schematic; per-galaxy profiles (see \S6) need not exhibit clean alternation.}
	\label{fig:concept_schematic}
\end{figure}

We proceed as follows. Section 2 develops the inverse-GR construction from kinematics to effective stress--energy. Section 3 specifies the minimal response-sector closure. Sections 4--5 report results across 175 galaxies. Sections 6--8 state falsification criteria and the domain-partitioned validation ladder. Section 9 provides context (MOND/TeVeS, emergent/nonlocal gravity), strengths and limitations, and the role-versus-identity framing. Section 10 concludes.
